\section{ВВЕДЕНИЕ}

Новые терапевтические подходы к лечению сердечно-сосудистых заболеваний требуют глубокого понимания молекулярных механизмов, лежащих в основе обновления и репарации сердечной ткани. Несмотря на значительный прогресс в методах профилактики и лечения различных патологий, связанных с разрушением тканей сердца, сердце остаётся органом с ограниченной способностью к самовосстановлению, что является одной из причин высокой летальности и ограниченной реабилитации пациентов с проблемами сердечно-сосудистой системы \cite{martin2023new}.

Изучение на молекулярном уровне процессов обновления и репарации повреждённых ткани приобретает особую актуальность для разработки новых стратегий лечения, профилактики и схем реабилитации при различных патологиях сердца. Последние исследования показывают важную роль в процессах репарации миокарда сигнальных путей регуляции эндотелиальной NO-синтазы (eNOS, Reactome: R-HSA-202127) \cite{zou2019vegf, dominguez2018urocortin, tran2022endothelial} и активации гипоксически-индуцированного фактора 1a (HIF-1a, Reactome: R-HSA-1234174.5), необходимого для образования кровеносных сосудов из уже существующей сосудистой сети \cite{gonzalez2017hypoxia}. В процессах ангиогенеза, важных для восстановления и репарации сердечной ткани, также активно участвует сигнальный путь VEGF (Reactome: R-HSA-194138.4)\cite{kobayashi2017dynamics, ferraro2019pro}, а также сигнальные пути других цитокинов и факторов роста: фактора роста фибробластов 2 (FGF2, Reactome: R-HSA-190236.4), тромбоцитарный фактор роста (PDGF, Reactome: R-HSA-186797.6), инсулиновый фактор роста-1 (IGF-1, Reactome: R-HSA-2404192.5), интерлейкины 2, 6, 17 (IL2: R-HSA-451927.7, IL6: R-HSA-6783589.8, IL17: R-HSA-448424.8), фактор некроза опухолей (TNF-a, Reactome: R-HSA-75893.10) \cite{virag2007fgf2, ono1997hgf,tomita2003hgf, zhao2011pdgf, dobrucki2011igf1,iwakura2006estradiol, barnabas2013estrogen,bouchentouf2011il2, mayfield2017il6, kishore2011tnf,xia2007mcp1}.

В последние годы исследовательский интерес всё чаще привлекают малые некодирующие регуляторные РНК (микроРНК). Эти молекулы участвуют в координации каскадов сигнальных путей, в т.ч. задействованных в пролиферации и дифференцировки клеток, апоптозе и регуляции воспалительных процессов, играющих важную роль в регенерации сердечной ткани \cite{wang2021myocardial, wang2021non}. Наши последние наработки в этой области позволили связать функцию микроРНК с топологическими характеристиками сетей белок-белковых взаимодействий продуктов их генов-мишеней \cite{osmak2020mirna,osmak2024MolBiol,pisklova2024unveiling}. Результаты наших исследований подводят к тому, что топологические характеристики сетей межмолекулярных взаимодействий могут быть уникальным функциональными сигнатурами множества генов, образующих эти сети. В настоящей работе, мы анализируя топологию сетей сигнальных путей постараемся выбрать наиболее топологически комплиментарные им сети генов-мишеней микроРНК. Таким образом, зная сигнальные пути участвующие в процессах обновления и репарации сердечной ткани, мы сможем очертить множество микроРНК, которое наиболее вероятно будет участвовать в регуляции этих процессов.